\documentclass[a5paper,10pt]{ltjsarticle}
\title{PDF の開き方入門}\author{Book Viewer 見本}\date{}
\begin{document}\maketitle
\section{開き方の設定}
これらはカタログ辞書の /PageLayout と /ViewerPreferences の /Direction に書かれます。対応したビューアーは、ファイルを開いたときにこの設定に従って表示します。これらはカタログ辞書の /PageLayout と /ViewerPreferences の /Direction に書かれます。対応したビューアーは、ファイルを開いたときにこの設定に従って表示します。

表示中の PDF が更新されたときは、書き込みが終わるのを待ってから読み込み直します。末尾の \%\%EOF を確かめるので、転送途中の中途半端なファイルは読みません。表示中の PDF が更新されたときは、書き込みが終わるのを待ってから読み込み直します。末尾の \%\%EOF を確かめるので、転送途中の中途半端なファイルは読みません。

表示中の PDF が更新されたときは、書き込みが終わるのを待ってから読み込み直します。末尾の \%\%EOF を確かめるので、転送途中の中途半端なファイルは読みません。PDF ファイルは、文書のプロパティとして「開き方」を持つことができます。ページレイアウト（単一ページ・見開き）と綴じ方（左・右）がその代表です。

これらはカタログ辞書の /PageLayout と /ViewerPreferences の /Direction に書かれます。対応したビューアーは、ファイルを開いたときにこの設定に従って表示します。表示中の PDF が更新されたときは、書き込みが終わるのを待ってから読み込み直します。末尾の \%\%EOF を確かめるので、転送途中の中途半端なファイルは読みません。

表示中の PDF が更新されたときは、書き込みが終わるのを待ってから読み込み直します。末尾の \%\%EOF を確かめるので、転送途中の中途半端なファイルは読みません。書き込みは一時ファイルに保存し、読み戻して検証してから元のファイルと差し替えます。途中で失敗しても、元のファイルは壊れません。

これらはカタログ辞書の /PageLayout と /ViewerPreferences の /Direction に書かれます。対応したビューアーは、ファイルを開いたときにこの設定に従って表示します。表示中の PDF が更新されたときは、書き込みが終わるのを待ってから読み込み直します。末尾の \%\%EOF を確かめるので、転送途中の中途半端なファイルは読みません。

\begin{center}\begin{tabular}{lll}\hline 項目 & 値 & 意味\\\hline
/PageLayout & /TwoPageRight & 見開き（表紙あり）\\ /Direction & /R2L & 右綴じ\\\hline\end{tabular}\end{center}
\section{書き込みの手順}
書き込みは一時ファイルに保存し、読み戻して検証してから元のファイルと差し替えます。途中で失敗しても、元のファイルは壊れません。表示中の PDF が更新されたときは、書き込みが終わるのを待ってから読み込み直します。末尾の \%\%EOF を確かめるので、転送途中の中途半端なファイルは読みません。

書き込みは一時ファイルに保存し、読み戻して検証してから元のファイルと差し替えます。途中で失敗しても、元のファイルは壊れません。表示中の PDF が更新されたときは、書き込みが終わるのを待ってから読み込み直します。末尾の \%\%EOF を確かめるので、転送途中の中途半端なファイルは読みません。

これらはカタログ辞書の /PageLayout と /ViewerPreferences の /Direction に書かれます。対応したビューアーは、ファイルを開いたときにこの設定に従って表示します。これらはカタログ辞書の /PageLayout と /ViewerPreferences の /Direction に書かれます。対応したビューアーは、ファイルを開いたときにこの設定に従って表示します。

PDF ファイルは、文書のプロパティとして「開き方」を持つことができます。ページレイアウト（単一ページ・見開き）と綴じ方（左・右）がその代表です。これらはカタログ辞書の /PageLayout と /ViewerPreferences の /Direction に書かれます。対応したビューアーは、ファイルを開いたときにこの設定に従って表示します。

これらはカタログ辞書の /PageLayout と /ViewerPreferences の /Direction に書かれます。対応したビューアーは、ファイルを開いたときにこの設定に従って表示します。これらはカタログ辞書の /PageLayout と /ViewerPreferences の /Direction に書かれます。対応したビューアーは、ファイルを開いたときにこの設定に従って表示します。

これらはカタログ辞書の /PageLayout と /ViewerPreferences の /Direction に書かれます。対応したビューアーは、ファイルを開いたときにこの設定に従って表示します。PDF ファイルは、文書のプロパティとして「開き方」を持つことができます。ページレイアウト（単一ページ・見開き）と綴じ方（左・右）がその代表です。

\begin{center}\begin{tabular}{lll}\hline 項目 & 値 & 意味\\\hline
/PageLayout & /TwoPageRight & 見開き（表紙あり）\\ /Direction & /R2L & 右綴じ\\\hline\end{tabular}\end{center}
\section{自動再読み込み}
表示中の PDF が更新されたときは、書き込みが終わるのを待ってから読み込み直します。末尾の \%\%EOF を確かめるので、転送途中の中途半端なファイルは読みません。これらはカタログ辞書の /PageLayout と /ViewerPreferences の /Direction に書かれます。対応したビューアーは、ファイルを開いたときにこの設定に従って表示します。

書き込みは一時ファイルに保存し、読み戻して検証してから元のファイルと差し替えます。途中で失敗しても、元のファイルは壊れません。書き込みは一時ファイルに保存し、読み戻して検証してから元のファイルと差し替えます。途中で失敗しても、元のファイルは壊れません。

PDF ファイルは、文書のプロパティとして「開き方」を持つことができます。ページレイアウト（単一ページ・見開き）と綴じ方（左・右）がその代表です。これらはカタログ辞書の /PageLayout と /ViewerPreferences の /Direction に書かれます。対応したビューアーは、ファイルを開いたときにこの設定に従って表示します。

表示中の PDF が更新されたときは、書き込みが終わるのを待ってから読み込み直します。末尾の \%\%EOF を確かめるので、転送途中の中途半端なファイルは読みません。書き込みは一時ファイルに保存し、読み戻して検証してから元のファイルと差し替えます。途中で失敗しても、元のファイルは壊れません。

書き込みは一時ファイルに保存し、読み戻して検証してから元のファイルと差し替えます。途中で失敗しても、元のファイルは壊れません。これらはカタログ辞書の /PageLayout と /ViewerPreferences の /Direction に書かれます。対応したビューアーは、ファイルを開いたときにこの設定に従って表示します。

PDF ファイルは、文書のプロパティとして「開き方」を持つことができます。ページレイアウト（単一ページ・見開き）と綴じ方（左・右）がその代表です。表示中の PDF が更新されたときは、書き込みが終わるのを待ってから読み込み直します。末尾の \%\%EOF を確かめるので、転送途中の中途半端なファイルは読みません。

\begin{center}\begin{tabular}{lll}\hline 項目 & 値 & 意味\\\hline
/PageLayout & /TwoPageRight & 見開き（表紙あり）\\ /Direction & /R2L & 右綴じ\\\hline\end{tabular}\end{center}
\section{空白ページの挿入}
表示中の PDF が更新されたときは、書き込みが終わるのを待ってから読み込み直します。末尾の \%\%EOF を確かめるので、転送途中の中途半端なファイルは読みません。表示中の PDF が更新されたときは、書き込みが終わるのを待ってから読み込み直します。末尾の \%\%EOF を確かめるので、転送途中の中途半端なファイルは読みません。

表示中の PDF が更新されたときは、書き込みが終わるのを待ってから読み込み直します。末尾の \%\%EOF を確かめるので、転送途中の中途半端なファイルは読みません。表示中の PDF が更新されたときは、書き込みが終わるのを待ってから読み込み直します。末尾の \%\%EOF を確かめるので、転送途中の中途半端なファイルは読みません。

PDF ファイルは、文書のプロパティとして「開き方」を持つことができます。ページレイアウト（単一ページ・見開き）と綴じ方（左・右）がその代表です。表示中の PDF が更新されたときは、書き込みが終わるのを待ってから読み込み直します。末尾の \%\%EOF を確かめるので、転送途中の中途半端なファイルは読みません。

表示中の PDF が更新されたときは、書き込みが終わるのを待ってから読み込み直します。末尾の \%\%EOF を確かめるので、転送途中の中途半端なファイルは読みません。PDF ファイルは、文書のプロパティとして「開き方」を持つことができます。ページレイアウト（単一ページ・見開き）と綴じ方（左・右）がその代表です。

これらはカタログ辞書の /PageLayout と /ViewerPreferences の /Direction に書かれます。対応したビューアーは、ファイルを開いたときにこの設定に従って表示します。PDF ファイルは、文書のプロパティとして「開き方」を持つことができます。ページレイアウト（単一ページ・見開き）と綴じ方（左・右）がその代表です。

これらはカタログ辞書の /PageLayout と /ViewerPreferences の /Direction に書かれます。対応したビューアーは、ファイルを開いたときにこの設定に従って表示します。表示中の PDF が更新されたときは、書き込みが終わるのを待ってから読み込み直します。末尾の \%\%EOF を確かめるので、転送途中の中途半端なファイルは読みません。

\begin{center}\begin{tabular}{lll}\hline 項目 & 値 & 意味\\\hline
/PageLayout & /TwoPageRight & 見開き（表紙あり）\\ /Direction & /R2L & 右綴じ\\\hline\end{tabular}\end{center}
\end{document}
